\section{Introduction}\label{sec:intro}

The Border Gateway Protocol (BGP) lets Autonomous Systems (ASes) exchange routes to IP prefixes.
An announcement alone does not show whether its origin AS may announce the prefix.
The Resource Public Key Infrastructure (RPKI) supplies this information through Route Origin Authorizations (ROAs).
Route origin validation (ROV) compares announcements with the ROAs and lets operators reject invalid routes~\cite{rfc6811}.
Earlier work shows that ROV limits how far invalid routes travel~\cite{testart2020filter,li2023rovista}.

Reach is only one part of the picture.
When an origin changes, the new route spreads over seconds to minutes, and the previous route can return later.
How fast the route spreads decides how much an operator or a detector sees early, and what is still to come.
Controlled experiments measure this for induced hijacks~\cite{sermpezis2018artemis}, but few studies follow it for the origin changes and leaks that occur in practice.
We measure the time course from public archives, for every RPKI state and for the return of the previous origin.

We ask two questions.
First, how do origin changes with different RPKI states spread, and how does the previous route return?
Second, how well do the first seconds of an event forecast the ASes that it reaches later, and which methods add to simple baselines?
Both answers depend on when the clock starts, which events a study selects, and what the collectors see.
If most of the final AS set is already visible early, persistence, the forecast that keeps the observed set, is hard to beat.
If the event keeps spreading, past events and the AS graph can help to find the ASes that join later.

\paragraph{Approach.}
We replay the archives of ten RIPE RIS and RouteViews collectors from January 2022 to August 2026 (\Cref{sec:data}).
The replay selections hold \num{66026} origin-change candidates and \num{17247} route-leak candidates.
A separate stratum of \num{13844} RPKI-invalid origin changes, drawn without labels, tests how catalogue selection shapes the results.

Historical ROAs classify the new and the previous origin of each event.
For each event we rebuild the set of ASes on observed paths over two hours.
This set is what the collectors see, not the complete set of affected networks~\cite{oliveira2009completeness}.
An invalid origin alone is not evidence of an attack, so we call these events origin changes, not hijacks.

For the forecasts, we use three setups of increasing rigour, and the first, on large leaks, is exploratory (\Cref{tab:forecast-setups}).
R1, a retrospective evaluation, covers selected leak candidates of late 2025, matched to Cloudflare Radar or high-volume.
It is descriptive, because we developed the methods while its period was open.
A prespecified replication on a new period, September to October 2026 (R2), repeats R1 with its arms, tests and outcome wording fixed in advance.
We operated its seal ourselves, no third party timestamped it, and an AI session took part in its decisions (\Cref{sec:genai}).

\paragraph{Contributions.}
Our four contributions show how the clock, the selection or the collectors shape early observations.
\begin{itemize}
\item \textbf{The clock and the selection decide class comparisons (\Cref{sec:prop}).} From the first observation of the new origin, invalid-origin events show the largest share of their two-hour AS set at 30\,s.
It is 0.61--0.74 in 2023--2026, against 0.32--0.42 for valid origins.
The new origin was visible more than a minute before the identified change for 54--72\,\% of valid events, but for only 6--16\,\% of invalid events.
For events first observed near the change, class differences in 2023--2026 shrink, but this rule keeps 84--94\,\% of invalid and 28--46\,\% of valid events.
In 2023--2025, a label-free sample shows that catalogue selection also raises the 30-s share of invalid events.
\item \textbf{The return differs by identifier gap (\Cref{sec:recovery}).} For pairs with an identifier gap of at least 15 minutes, the return shows more of its final set at 30\,s.
Its day-block intervals are above zero in 2022, 2023 and 2025.
Most of the \num{20256} matched pairs, 55--65\,\%, are flips less than one minute apart.
In these flips the observed return set is larger than the outward set.
An early end of the outward change, ROV or path exploration may explain it, and we did not test which.
\item \textbf{The collectors see a changing share of invalid routes (\Cref{sec:beacons}).} The relative reach of a controlled invalid beacon fell by about half between January and May 2024.
Two upstream changes in April coincide with the steps of this fall, on top of an earlier decline.
\item \textbf{Forecast comparisons depend on the selection and the setup (\Cref{sec:forecast}).} On large leaks selected by final size, graph models gain up to $+0.20$ Jaccard over persistence in exploratory tests, and history almost as much.
In R1, which differs in selection, training and calibration, the set forecast of SAGE with history equals persistence for 95\,\% of events at 30\,s.
A revision analysis finds it within 0.001 Jaccard of persistence by point estimate also on fast-growing events.
% round 6 (panel-r6 NEW-4(b) / X5; review-r6 "add by point estimate"): the slice intervals reach -0.0028 and +0.0014, and the R1 30-s
%   interval [+0.0003, +0.0014] excludes zero (sections/08-forecast.tex "Secondary metrics", tab:slice-persistence); "events that grow by issue time"
%   -> "events with fast early growth" (the slice: early growth >= 12 at 30 s, >= 22 at 60 s; 966 of 3,543 R1 events at 30 s).

Message passing and history each add about 0.06 average precision (AP) alone, and 0.01--0.02 given the other.
On the same candidates, we detect no AP difference between SAGE with history and online history.
Online history with its own, nearly complete candidate set reaches a higher AP, in a comparison of complete systems.
A prespecified replication on autumn 2026 (R2) repeats these Holm decisions, and its other rejections favour history in ranking and SAGE in Jaccard.
\end{itemize}

We release the code and the derived data (\Cref{sec:openscience}).
% src: see sections/00-abstract.tex; recovery strata from paper/generated/recovery_gap.tex; R1 from paper/generated/r1-comparisons.json.
% src: R1 two-by-two at 30 / 60 s (tab:r1, r1-comparisons.json): history alone C4 +0.063 / +0.059; message passing alone C1 + C4 - C3 = +0.068 / +0.063
%   (r1-arm-levels.json sage_off - none_off 0.0679 / 0.0635); given the other: C1 +0.018 / +0.017, C3 +0.014 / +0.013.
%   month boundary: month_boundary.json R1 C2@30 no_day1 +0.0066 [+0.0008, +0.0118] (identifier date); R2 test events 1,289 (tab:forecast-setups).
%   return hypotheses: sections/06-recovery.tex "By identifier gap" (round-3 review NEW-7: truncation, ROV, path exploration).
% src: near-onset retention 84--94 % invalid, 28--46 % valid: paper/generated/onset_sensitivity_ci.json share_near, 2022--2026 (sections/05-propagation.tex);
%   R2 repeats the Holm decision of every common AP contrast: newperiod-table.tex vs r1-comparisons.json (C2 at 60 s changes sign, +0.0022 to -0.0042).
% src (round 4, 2026-10-31; panel-r4 T2 / NEW-1, NEW-2, NEW-5, NEW-12, Q24): gap >= 15 min, day-block intervals of paired Delta J above zero in
%   2022 [+0.093, +0.171], 2023 [+0.060, +0.141], 2025 [+0.035, +0.145], including zero in 2024 and 2026 (paper/generated/recovery_dayblock.tex);
%   heading "shapes" -> "differs by" (causal wording, NEW-12). Set forecast = persistence: set_errors.json R1.30.sage_on.adds_nothing_share 0.9463.
%   R2 P@10: facts.json p_at_10 C2@30 -0.0239, Holm 0.0024 (tab:secondary). The first-of-month sentence (+0.007, post hoc) left the bullet (NEW-1);
%   "ranks better" is kept for Holm rejections of ranking quality, and the native comparison now reads "reaches a higher AP" (NEW-5).
% Round 5 (2026-11-01; panel-r5 NEW-2, NEW-3, NEW-5, NEW-17; review-r5 items 2, 3, 6):
%   "but not for the origin changes and leaks" -> "few studies follow it" (prior-work claim narrowed; Argus, Shi et al. IMC 2012, not in refs.bib).
%   Return bullet: the larger return set is descriptive; the three hypotheses are untested (sections/06-recovery.tex).
%   Bullet 4: the slow regime and R1 also differ in training population (large leaks vs all selected candidates) and threshold calibration
%   (slow regime: threshold tuned on the training split of large leaks, sections/08-forecast.tex "Methods"; R1: validation mean Jaccard over all
%   events, tab:recipe "Learned set forecast"), so the contrast is not a within-pipeline estimate of the selection (review-r5 W2).
% src (round 5, panel-r5 A-R2; revision analysis, confirmatory false): paper/generated/slice_vs_persistence.json results.{R1,R2}.{30,60}
%   .arms.sage_on.minus_persistence.difference: R1 +0.000828 [+0.000255, +0.001387] (30 s), +0.0000387 (60 s); R2 -0.000837 [-0.002793, +0.000727]
%   (30 s), -0.0000119 (60 s); all |x| < 0.001. Slice: early growth >= 12 at 30 s, >= 22 at 60 s (definitions.slice; tab:recipe). The slice is
%   the look-ahead-free analogue of the slow regime: an issue-time selection of growing events alone does not move the set forecast off persistence.
%   The bullet is split into a set-forecast and a ranking paragraph (six sentences at most per paragraph).
% src: R2 rejections by direction: sections/00-abstract.tex round-5 src comment (facts.json holm_rejects): AP C5@60 and P@10 C2@30, C5@60 and the
%   three native-family AP contrasts favour online history; full-target and slice Jaccard C2@30, C2@60, C5@60 favour SAGE.
